% 模板 5/6/7/8 合集 — 统一精致配色（蓝色主节点+teal子节点+微阴影）
% 编译命令：xelatex demo_all_templates.tex
\documentclass[a4paper,12pt]{article}
\usepackage[margin=1.5cm]{geometry}
\usepackage{ctex}
\usepackage{tikz}
\usepackage{float}
\usetikzlibrary{arrows.meta, positioning, calc, shapes.geometric, shadows}

\begin{document}

% ==================== 模板 5 ====================
\section*{模板 5：算法流程图（带判断分支）}
\begin{figure}[H]
\centering
\begin{tikzpicture}[
    node distance=0.9cm,
    process/.style={rectangle, rounded corners=3pt,
        minimum width=4cm, minimum height=0.7cm, align=center,
        draw=blue!80, line width=0.7pt, fill=blue!6,
        font=\small, drop shadow={opacity=0.15, shadow xshift=0.5pt, shadow yshift=-0.5pt}},
    decision/.style={diamond, aspect=2.5,
        minimum width=2cm, align=center,
        draw=orange!70, line width=0.7pt, fill=orange!6,
        font=\small, inner sep=1pt,
        drop shadow={opacity=0.12, shadow xshift=0.5pt, shadow yshift=-0.5pt}},
    io/.style={rectangle, rounded corners=3pt,
        minimum width=3.5cm, minimum height=0.7cm, align=center,
        draw=teal!70, line width=0.5pt, fill=white, font=\small},
    arrow/.style={-stealth, line width=0.7pt, color=gray!60},
    yesno/.style={font=\footnotesize, color=gray!50},
]
\node[io] (start) {输入数据 $D$};
\node[process, below=of start] (init) {初始化参数 $\theta_0$};
\node[process, below=of init] (compute) {计算目标函数 $f(\theta)$};
\node[process, below=of compute] (update) {更新参数 $\theta \leftarrow \theta - \alpha\nabla f$};
\node[decision, below=of update] (conv) {收敛?};
\node[io, below=of conv] (output) {输出最优解 $\theta^*$};
\draw[arrow] (start) -- (init); \draw[arrow] (init) -- (compute);
\draw[arrow] (compute) -- (update); \draw[arrow] (update) -- (conv);
\draw[arrow] (conv) -- node[yesno, right] {是} (output);
\draw[arrow] (conv.west) -- ++(-1.8,0) node[yesno, above] {否} |- (compute.west);
\end{tikzpicture}
\caption{模板 5：算法流程图}
\end{figure}

\newpage

% ==================== 模板 6 ====================
\section*{模板 6：数据处理 Pipeline（横向多阶段）}
\begin{figure}[H]
\centering
\begin{tikzpicture}[
    node distance=0.3cm,
    stage/.style={rectangle, rounded corners=3pt,
        minimum width=2.4cm, minimum height=2.2cm, align=center,
        draw=blue!80, line width=0.7pt, fill=blue!6,
        font=\small,
        drop shadow={opacity=0.15, shadow xshift=0.5pt, shadow yshift=-0.5pt}},
    detail/.style={rectangle, rounded corners=2pt,
        minimum width=2.2cm, minimum height=0.5cm, align=center,
        draw=teal!70, line width=0.5pt, fill=white,
        font=\tiny, text width=2cm},
    arrow/.style={-stealth, line width=1.2pt, color=blue!70},
]
\node[stage] (raw) {\textbf{原始数据}\\[2pt]\footnotesize 多源采集};
\node[stage, right=1.2cm of raw] (clean) {\textbf{数据清洗}\\[2pt]\footnotesize 缺失值/异常值};
\node[stage, right=1.2cm of clean] (feat) {\textbf{特征工程}\\[2pt]\footnotesize 变量构建};
\node[stage, right=1.2cm of feat] (model) {\textbf{模型训练}\\[2pt]\footnotesize 参数优化};
\node[stage, right=1.2cm of model] (eval) {\textbf{评估验证}\\[2pt]\footnotesize 交叉验证};
\node[detail, below=0.4cm of raw] {CSV/API/\\数据库};
\node[detail, below=0.4cm of clean] {插值/IQR/\\标准化};
\node[detail, below=0.4cm of feat] {PCA/交互项/\\时序特征};
\node[detail, below=0.4cm of model] {XGBoost/\\DNN/SVM};
\node[detail, below=0.4cm of eval] {RMSE/AUC/\\$R^2$};
\draw[arrow] (raw) -- (clean); \draw[arrow] (clean) -- (feat);
\draw[arrow] (feat) -- (model); \draw[arrow] (model) -- (eval);
% 反馈箭头走上方
\draw[arrow, dashed, color=red!40] (eval.north) -- ++(0,0.8) -| (feat.north)
    node[pos=0.25, above, font=\tiny, color=red!50] {特征调优};
\end{tikzpicture}
\caption{模板 6：数据处理 Pipeline}
\end{figure}

\newpage

% ==================== 模板 7 ====================
\section*{模板 7：变量关系路径图（中介效应）}
\begin{figure}[H]
\centering
\begin{tikzpicture}[
    node distance=2cm and 3cm,
    var/.style={rectangle, rounded corners=3pt,
        minimum width=3.2cm, minimum height=1cm, align=center,
        draw=blue!80, line width=0.7pt, fill=blue!6,
        font=\small,
        drop shadow={opacity=0.15, shadow xshift=0.5pt, shadow yshift=-0.5pt}},
    ctrl/.style={rectangle, rounded corners=2pt,
        minimum width=2.8cm, minimum height=0.7cm, align=center,
        draw=teal!70, line width=0.5pt, fill=white, font=\footnotesize},
    arrow/.style={-stealth, line width=0.8pt},
    coef/.style={font=\footnotesize, fill=white, inner sep=2pt},
]
\node[var] (x) {\textbf{自变量}\\数字化转型};
\node[var, above right=1.5cm and 4cm of x] (m) {\textbf{中介变量}\\创新能力};
\node[var, below right=1.5cm and 4cm of x] (y) {\textbf{因变量}\\企业绩效};
\draw[arrow, color=blue!60] (x) -- node[coef, below] {$c'$ (直接效应)} (y);
\draw[arrow, color=orange!60] (x) -- node[coef, above left] {$a$ (H1)} (m);
\draw[arrow, color=red!60] (m) -- node[coef, above right] {$b$ (H2)} (y);
\node[ctrl, below=1.5cm of y] (ctrl) {控制变量\\企业规模/行业/年份};
\draw[-stealth, dashed, color=gray!40, line width=0.5pt] (ctrl) -- (y);
\end{tikzpicture}
\caption{模板 7：变量关系路径图}
\end{figure}

\newpage

% ==================== 模板 8 ====================
\section*{模板 8：单问题求解流程图（带分支+判断+并行）}
\begin{figure}[H]
\centering
\begin{tikzpicture}[
    node distance=0.8cm and 1.2cm,
    step/.style={rectangle, rounded corners=3pt,
        minimum width=3.8cm, minimum height=0.7cm, align=center,
        draw=blue!80, line width=0.7pt, fill=blue!6,
        font=\small,
        drop shadow={opacity=0.15, shadow xshift=0.5pt, shadow yshift=-0.5pt}},
    substep/.style={rectangle, rounded corners=2pt,
        minimum width=2.8cm, minimum height=0.6cm, align=center,
        draw=teal!70, line width=0.5pt, fill=white, font=\footnotesize},
    decision/.style={diamond, aspect=2.8,
        minimum width=1.5cm, align=center,
        draw=orange!70, line width=0.7pt, fill=orange!6,
        font=\small, inner sep=1pt,
        drop shadow={opacity=0.12, shadow xshift=0.5pt, shadow yshift=-0.5pt}},
    dashbox/.style={rectangle, rounded corners=4pt,
        draw=gray!40, dashed, line width=0.7pt, fill=gray!2},
    arrow/.style={-stealth, line width=0.5pt, color=gray!55},
    bigarrow/.style={-stealth, line width=1.4pt, color=blue!70},
    yesno/.style={font=\tiny, color=gray!50},
    label/.style={font=\small\bfseries, color=black!70},
]

% 阶段一：数据准备
\node[dashbox, minimum width=12cm, minimum height=3.5cm] (box1) at (0, 0) {};
\node[label, anchor=north west] at (box1.north west) {\scriptsize 数据准备};
\node[step] (input) at (0, 1.0) {读取附件数据};
\node[step, below=0.6cm of input] (eda) {数据探索与可视化};
\node[decision, below=0.7cm of eda] (missing) {有缺失值?};
\node[substep, right=1.8cm of missing] (fill) {插值/删除处理};
\draw[arrow] (input) -- (eda); \draw[arrow] (eda) -- (missing);
\draw[arrow] (missing) -- node[yesno, above] {是} (fill);
\draw[arrow] (missing) -- node[yesno, right] {否} ++(0,-1.0);

\draw[bigarrow] (0, -2.5) -- (0, -3.2);

% 阶段二：建模（并行）
\node[dashbox, minimum width=12cm, minimum height=3cm] (box2) at (0, -4.8) {};
\node[label, anchor=north west] at (box2.north west) {\scriptsize 模型构建};
\node[step] (formulate) at (0, -3.6) {建立数学模型};
\node[substep] (method1) at (-3, -4.8) {方法A：精确求解};
\node[substep] (method2) at (3, -4.8) {方法B：启发式$^{\bigstar}$};
\draw[arrow] (formulate.south) -- ++(0,-0.3) -| (method1.north);
\draw[arrow] (formulate.south) -- ++(0,-0.3) -| (method2.north);
\node[step, below=1.5cm of formulate] (compare) {方法对比与选优};
\draw[arrow] (method1.south) |- (compare.west);
\draw[arrow] (method2.south) |- (compare.east);

\draw[bigarrow] (0, -6.5) -- (0, -7.2);

% 阶段三：验证
\node[dashbox, minimum width=12cm, minimum height=2.2cm] (box3) at (0, -8.4) {};
\node[label, anchor=north west] at (box3.north west) {\scriptsize 结果验证};
\node[step] (verify) at (-2.5, -8.0) {结果验证与分析};
\node[step] (sense) at (2.5, -8.0) {灵敏度分析};
\node[step] (output) at (0, -9.1) {输出最终方案};
\draw[arrow] (verify) -- (output); \draw[arrow] (sense) -- (output);

\end{tikzpicture}
\caption{模板 8：单问题求解流程图}
\end{figure}

\end{document}
